% !TEX root = ../main.tex
\lecture{1}{2026-09-13}{Social Conscience}

\section{Social being and forms of thought}

\begin{quote}
   Es ist nicht das Bewusstsein der Menschen, das ihr Sein, sondern umgekehrt, ihr gesellschaftliches Sein, das ihr Bewusstsein bestimmt. 
   \author{K. Marx},
   Zur Kritik der politischen konomie
\end{quote}

Above is one of the quotes that is most misunderstood by liberals and conservatives worldwide. Taken in its vulgar sense, the phrase might sound  
like a mechanicist, determinist law which subjugates the "free-will" of individuals to their material and social being. The misunderstanding lies in
the verb \textit{bestimmt}, which is often interpreted as an affirmation of direct causality, and thus provides a base for interpretations in which a capitalist society
would inherently produce its members think as capitalists, or believe in absolute equality and freedom without room for critical thought or individual will.

If we accept this understanding of the sentence and return to the broader context of Marx' work, then its contradiction becomes blatant, for if he truly meant 
to describe consciousness as immediately caused by the social being, then how would he have been able to develop
such a critical analysis of the society he found himself in? Did he consider himself an exception to the rule? What about engels and all the other thousands of less known communist activists, who dreamed of overcoming said society?
From this basic fact it becomes apparent that this interpretation is absolutely flawed and a more rigorous look into his work make it an impossibility.

Furthermore it is a mistake to take that quote as a hard truth or "method of analysis" by itself, while it is nothing but the result of a process we haven't completed ourselves.
 It does not provide us with any relevant knowledge nor answer the questions of which and why specific forms of conscience/thought tend to appear in capitalist society, or in any other given form of society. 
 The only way, and this is the materialist assumption, to gain knowledge is to look into each of the individual societies we seek to investigate and the  internal connections of their determining factors and find out for ourselves if its forms of thought correlate in one way or another with the maintenance or criticism of said conditions.
The search for universal law, dogma or axiom of the general course of society and its forms of thought is either trivial or purely metaphysical, so if we were to reach the conclusion that in every society there must be working people producing and this production has to be distributed in any way it is as empty of content as the philosophies of history
,which part from an assumed benevolent or egoistic human nature or absolute concepts/ideas and deduce from them the whole course of human nature.
Although it is to say that even if mistaken in principles such methods have had in their course of analysis some important discoveries or indirectly contributed to a broader understanding of reality. So per example the ancient form of astrology or alchemy contributed in some way or another, in fragments of correct knowledge or in form of research to navigation, modern astronomy and chemistry.
 Kant, who despite being a philosopher of metaphysical tradition, came to correctly deduce the material and not miraculous origin of galaxies and our solar system. 
 Hegel whose understanding of internal contradictions and form of presentation heavily aligns with the appearance and development of capitalist social forms such as the unfolding of the primal contradiction of private ownership and social distribution.


